\documentclass[11pt, a4paper]{article}

\usepackage[utf8]{inputenc}
\usepackage[T1]{fontenc}
\usepackage[french]{babel}
\usepackage{amsmath, amssymb, amsfonts}
\usepackage{geometry}
\usepackage{graphicx}
\usepackage{hyperref}
\usepackage{cleveref}
\usepackage{booktabs}
\usepackage{xcolor}
\usepackage{float}
\usepackage{tikz}
\usepackage{siunitx}
\usetikzlibrary{shapes.geometric, arrows.meta, positioning, calc}
\usepackage{enumitem}
\usepackage{changepage}

\geometry{margin=2.5cm}

\title{\textbf{Modélisation et Apprentissage pour la Simulation d’Éoliennes et leur Optimisation}}
\author{Adam Rozentalis}
\date{\today}


\begin{document}

\maketitle

\section{Introduction et Objectifs}
Le calcul des efforts aérodynamiques sur les pales d'une éolienne est traditionnellement soumis à un compromis. D'un côté, la méthode \textbf{BEM (Blade Element Momentum)} offre une résolution extrêmement rapide, idéale pour l'optimisation, mais peine à capturer les effets tridimensionnels complexes. De l'autre, les méthodes \textbf{Vortex} (sillage libre) offrent une excellente précision en modélisant les interactions de sillage, mais avec un coût de calcul important.\\

L'objectif de ce stage est de développer un algorithme de \textbf{Deep Learning} capable d'inférer les efforts aérodynamiques avec la \textbf{précision de la méthode Vortex}, tout en conservant une \textbf{complexité algorithmique comparable à celle de la BEM}. L'avantage des réseaux de neurones est qu'ils sont différentiables ce qui peut être très utile pour l'optimisation du design des éoliennes.

\section{Architecture des Modèles et Espaces des Paramètres}

Pour que le réseau de neurones puisse apprendre efficacement la physique du problème, il est nécessaire de bien définir, simplifier et structurer les espaces d'entrée ($\mathcal{X}$) et de sortie ($\mathcal{Y}$). La \Cref{fig schema BEM et Vortex} compare le fonctionnement de la méthode BEM à la méthode Vortex. Cette comparaison conduit naturellement à envisager trois grandes familles de stratégies : 
\begin{itemize}
	\item la discrétisation
	\item la prédiction résiduelle 
	\item les variables intermédiaires
\end{itemize}

\subsection{L'Approche Locale (BEM) versus l'Approche Globale (Vortex)}

L'approche locale schématisée dans la \Cref{fig modele_local} reproduit l'hypothèse d'indépendance radiale de la BEM : chaque section de la pale est traitée isolément. L'espace d'entrée $\mathcal{X}_{com}\times \mathcal{X}_{loc}$ contient l'information d'un seul tronçon. \\

L'espace $\mathcal{X}_{com}\times\mathcal{X}_{loc}$ comprend 1 variable discrète ($B \in \mathbb{N}^*$), 11 variables continues ($r$ , $\theta$, $R$, $R_{hub}$, $c$, $\beta$, $\theta_{pitch}$, $\rho$, $\Omega$, $U_\infty$, $\gamma$) et 2 variables fonctionnelles ($C_L$, $C_D)$. \\

On précise que la vitesse de rotation $\Omega$ et la vitesse du vent $U_\infty$ qui y figurent sont physiquement couplées par le \textit{Tip Speed Ratio} (TSR), noté $\lambda = (\Omega R)/U_\infty$. Il est donc possible de substituer l'un de ces deux paramètres par $\lambda$. \\

En ce qui concerne la sortie $\mathcal{Y}$, $F_n$ et $F_t$ désigne les forces calculées par la méthode Vortex dans le repère local de la pale. On aurait pu choisir comme target la poussée axiale ($T$) et le couple ($Q$) définit dans le repère du rotor. Le passage d'un repère à l'autre se fait par une simple rotation impliquant l'angle de twist $\beta$ et l'angle de pitch $\theta_{pitch}$ :
    $$ \begin{pmatrix} T \\ Q/r \end{pmatrix} = \begin{pmatrix} \cos(\theta_{pitch} + \beta) & -\sin(\theta_{pitch} + \beta) \\ \sin(\theta_{pitch} + \beta) & \cos(\theta_{pitch} + \beta) \end{pmatrix} \begin{pmatrix} F_n \\ F_t \end{pmatrix}.$$


\begin{adjustwidth}{-1cm}{-1cm}
	\begin{figure}
	\begin{tikzpicture}[
		font=\large,
		node distance=1.5cm and 1.5cm,
		bubble/.style={
			draw, 
			rectangle, 
			rounded corners=15pt, % Rendu style "bulle" plus adapté aux longues listes
			align=left, 
			fill=blue!5,  
			inner sep=10pt,
			text width=6.3cm 
		},
		block/.style={
			draw, 
			rectangle, 
			fill=orange!20, 
			minimum width=2.5cm, 
			minimum height=1.5cm, 
			font=\bfseries\Large,
			rounded corners=2pt
		},
		arrow/.style={
			-{Stealth[scale=1.2]}, 
			line width=1pt, 
			draw=gray!80
		}
		] 
		
		\node[bubble] (Xloc) {
			$\mathcal{X}_{loc}$ : 
			\begin{itemize}[label=\textbullet]
				\item Position radiale $r \in [R_{hub},R]$
				\item Position azimutale $\theta \in [0,2\pi]$
				\item Corde $c \in \mathbb{R}^+$
				\item Twist $\beta \in [0,2\pi]$
				\item Coef lift $C_L \in \mathcal{C}([0,2\pi],\mathbb{R})$
				\item Coef drag $C_D \in \mathcal{C}([0,2\pi],\mathbb{R})$
			\end{itemize}
		};
		
		% Bloc X_com (haut)
		\node[bubble,below=2.9cm of Xloc] (Xcom) {
			$\mathcal{X}_{Com}$ : 
			\begin{itemize}[label=\textbullet]
				\item Nombre de pales $B \in \mathbb{N}^*$
				\item Rayon total $R \in \mathbb{R}^+$
				\item Rayon moyeu $R_{hub} \in \mathbb{R}^+$
				\item Pitch $\theta_{pitch} \in [0,2\pi]$
				\item Densité de l'air $\rho \in \mathbb{R}^+$
				\item Vitesse de rotation $\Omega \in \mathbb{R}$
				\item Vitesse freestream $U_{\infty} \in \mathbb{R}$
				\item Yaw $\gamma \in [-\pi/4, \pi/4]$
			\end{itemize}
		};
		
		% Bloc X_G (bas)
		\node[bubble, below=2.5cm of Xcom] (XG) {
			$\mathcal{X}_{G}$ : 
			\begin{itemize}[label=\textbullet]
				\item Positions radiales $r \in [R_{hub},R]^n$
				\item Positions azimutales $\theta \in [0,2\pi]^p$
				\item Corde $c \in \mathbb{R}^n$
				\item Twist $\beta \in [0,2\pi]^n$
				\item Coef lift $C_L \in \mathcal{C}([0,2\pi],\mathbb{R}^n)$
				\item Coef drag $C_D \in \mathcal{C}([0,2\pi],\mathbb{R}^n)$
			\end{itemize}
		};
		
		
		\node[block, right=0.5cm of Xcom, yshift=2.2cm] (BEM) {BEM};
		\node[block, right=0.5cm of Xcom,yshift=-2.2cm] (VORTEX) {VORTEX};
		
		
		% Sortie
		\node[bubble, fill=violet!15,right=2.2cm of XG, yshift=2.5cm,text width=4.5cm] (V) {
			$\mathcal{V}_G$ :
			\begin{itemize}[label=\textbullet]
				\item Vitesse effective $U_{eff} \in \mathcal{M}_{n,p}(\mathbb{R})$
				\item Angle d'attaque $\alpha \in \mathcal{M}_{n,p}(\mathbb{R})$
			\end{itemize}
		};
		
		\node[bubble,fill=violet!15, right=2.2cm of Xloc,yshift=-2.5cm, text width=4.5cm] (B) {
			$\mathcal{B}$ :
			\begin{itemize}[label=\textbullet]
				\item Induction axiale $a \in \mathbb{R}$
				\item Angle de flux $\varphi \in \mathbb{R}$
			\end{itemize}
		};
		
		\node[block, above=2cm of B] (Form1) {Formules};
		\node[block, below=2cm of V] (Form2) {Formules};
		
		% Bulle de sortie Y
		\node[bubble, fill=green!15,right=2cm of BEM, text width=4.5 cm] (Yloc) {
			$\mathcal{Y}_{loc} = \mathbb{R}^2$ : 
			\begin{itemize}[label=\textbullet]
				\item Force normale $F^{BEM}_n \in \mathbb{R}$
				\item Force tangentielle $F^{BEM}_t \in \mathbb{R}$
			\end{itemize}
		};
		
		\node[bubble, fill=green!15,right=2cm of VORTEX, text width=4.5 cm] (YG) {
			$\mathcal{Y}_G = \mathcal{M}_{n,p}(\mathbb{R})^2$ : 
			\begin{itemize}[label=\textbullet]
				\item Force normale $F^{Vortex}_n \in \mathcal{M}_{n,p}(\mathbb{R})$
				\item Force tangentielle $F^{Vortex}_t \in \mathcal{M}_{n,p}(\mathbb{R})$
			\end{itemize}
		};
		
		% Flèches
		\draw[arrow] (Xcom.east) -| (BEM.south);
		\draw[arrow] (Xcom.east) -| (VORTEX.north);
		\draw[arrow] (XG) -| (VORTEX);
		\draw[arrow] (Xloc) -| (BEM);
		
		% Sortie
		\draw[arrow] (BEM) -| (B);
		\draw[arrow] (VORTEX) -| (V);
		\draw[arrow] (V) -- (Form2);
		\draw[arrow] (B) -- (Form1);
		\draw[arrow] (Form1) -| (Yloc);
		\draw[arrow] (Form2) -| (YG);
		
	\end{tikzpicture}
\caption{Schéma des méthodes BEM et de Vortex pour le calcul des efforts aérodynamiques.}	
\label{fig schema BEM et Vortex}
\end{figure}
\end{adjustwidth}
	


De plus, il est important de noter que la gestion des fonctions aérodynamiques $C_L$ et $C_D$ pose un défi dimensionnel. On peut les évaluer sur une grille fixe de $m$ angles d'attaque, ou les décomposer sur une base de fonctions (ex : B-splines ou Fourier) pour réduire drastiquement cette dimension $m$.
De plus, dans la réalité physique, $C_L$ et $C_D$ ne dépendent pas uniquement de l'angle d'attaque $\alpha$, mais également du nombre de Reynolds ($Re$) qui pourrait donc constituer une variable d'entrée. Dans la suite, nous considérerons que l'information d'une courbe polaire est compressée dans un vecteur de taille $m$. \\

La dimension totale de l'espace d'entrée dans le modèle local est donc :
$$ d_{loc} = 12 + 2m. $$
Ce modèle cartographie $\mathcal{X}_{com}\times\mathcal{X}_{loc} \to \mathbb{R}^2$ pour calculer les efforts locaux $(F_n, F_t)$. \\

\begin{figure}[H]
	\centering
	\begin{tikzpicture}[
		node distance=1.5cm and 1cm,
		bubble/.style={draw, rectangle, rounded corners=10pt, align=center, fill=blue!5, inner sep=8pt, minimum height=1cm},
		block/.style={draw, rectangle, fill=orange!20, minimum width=2.5cm, minimum height=1cm, font=\bfseries, align=center},
		skynet/.style={draw, rectangle, fill=red!20, minimum width=2.5cm, minimum height=1cm, font=\bfseries\large, rounded corners=3pt},
		arrow/.style={-{Stealth[scale=1.2]}, line width=1pt, draw=gray!80}
		]
		
		% Noeuds
		\node[bubble] (X) {$x \in \mathcal{X}_{com}\times\mathcal{X}_{loc}$};
		\node[skynet, right=of X] (SkyNet) {SkyNet};
		\node[bubble, fill=green!15,right=of SkyNet] (Y) {$(\hat{F}_n, \hat{F}_t) \in \hat{\mathcal{Y}}_{loc}=\mathbb{R}^2$};
		
		% Flèches
		\draw[arrow] (X) -- (SkyNet);
		\draw[arrow] (SkyNet) -- (Y);
		
	\end{tikzpicture}
	\caption{Schéma du modèle local}
	\label{fig modele_local}
\end{figure}

Pour dépasser les limites de la BEM et saisir les interactions tridimensionnelles, l'approche globale schématisée dans la \Cref{fig modele_global} évalue la pale entière d'un seul coup pour $n \in \mathbb{N}^\star$ points radiaux (discrétisation spatiale) et $p$ points azimutaux (discrétisation temporelle) à l'instar de la méthode Vortex. A priori, le choix de la discrétisation n'influe pas sur la réalité physique, c'est une considération purement numérique. Par conséquent, on pourra se fixer des grilles de discrétisation radiale et azimutale.
Cette méthode se décline en trois sous variantes selon que l'on fixe les deux grilles ou une seule tout en faisant varier localement le paramètre restant. \\

\begin{figure}[H]
	\centering
	\begin{tikzpicture}[
		node distance=1.5cm and 1cm,
		bubble/.style={draw, rectangle, rounded corners=10pt, align=center, fill=blue!5, inner sep=8pt, minimum height=1cm},
		block/.style={draw, rectangle, fill=orange!20, minimum width=2.5cm, minimum height=1cm, font=\bfseries, align=center},
		skynet/.style={draw, rectangle, fill=red!20, minimum width=2.5cm, minimum height=1cm, font=\bfseries\large, rounded corners=3pt},
		arrow/.style={-{Stealth[scale=1.2]}, line width=1pt, draw=gray!80}
		]
		
		% Noeuds
		\node[bubble] (X) {$x \in \mathcal{X}_{com}\times\mathcal{X}_{G}$};
		\node[skynet, right=of X] (SkyNet) {SkyNet};
		\node[bubble, fill=green!15,right=of SkyNet] (Y) {$(\hat{F}_n, \hat{F}_t) \in \hat{\mathcal{Y}}_{G}$};
		
		% Flèches
		\draw[arrow] (X) -- (SkyNet);
		\draw[arrow] (SkyNet) -- (Y);
		
	\end{tikzpicture}
	\caption{Schéma du modèle global}
	\label{fig modele_global}
\end{figure}

\begin{figure}
\begin{tikzpicture}[
	font=\large,
	node distance=2cm and 2.5cm,
	bubble/.style={
		draw,
		rectangle,
		rounded corners=10pt,
		align=left,
		fill=blue!5,
		inner sep=8pt,
		text width=7cm
	},
	arrow/.style={
		-{Stealth[scale=1.2]},
		line width=1pt
	}
	]
	
	% --- Ligne 1 : Grt (complet r,theta) ---
	\node[bubble] (XGrt) { 
		$\mathcal{X}_{Grt}$ : grille de taille $(n,p)$ fixée
		\begin{itemize}[label=\textbullet]
			\item $c \in \mathbb{R}^n$
			\item $\beta \in [0,2\pi]^n$
			\item $C_L \in \mathcal{C}([0,2\pi],\mathbb{R}^n)$
			\item $C_D \in \mathcal{C}([0,2\pi],\mathbb{R}^n)$
		\end{itemize}
	};
	
	\node[bubble,fill=green!15, right=of XGrt] (YGrt) { 
		$\mathcal{Y}_{Grt} = \mathcal{M}_{n,p}(\mathbb{R})^2$
		\begin{itemize}[label=\textbullet]
			\item $F_n \in \mathcal{M}_{n,p}(\mathbb{R})$
			\item $F_t \in \mathcal{M}_{n,p}(\mathbb{R})$
		\end{itemize}
	};
	
	\draw[arrow] (XGrt) -- (YGrt);
	
	% --- Ligne 2 : Gr (radial) ---
	\node[bubble, below=of XGrt] (XGr) { 
		$\mathcal{X}_{Gr}$ : grille radiale de taille $n$ fixée
		\begin{itemize}[label=\textbullet]
			\item $\theta \in [0,2\pi]$
			\item $c \in \mathbb{R}^n$
			\item $\beta \in [0,2\pi]^n$
			\item $C_L \in \mathcal{C}([0,2\pi],\mathbb{R}^n)$
			\item $C_D \in \mathcal{C}([0,2\pi],\mathbb{R}^n)$
		\end{itemize}
	};
	
	\node[bubble,fill=green!15, right=of XGr] (YGr) { 
		$\mathcal{Y}_{Gr} = \mathbb{R}^{2n}$
		\begin{itemize}[label=\textbullet]
			\item $F_n \in \mathbb{R}^n$
			\item $F_t \in \mathbb{R}^n$
		\end{itemize}
	};
	
	\draw[arrow] (XGr) -- (YGr);
	
	% --- Ligne 3 : Gt (azimutal) ---
	\node[bubble, below=of XGr] (XGt) { 
		$\mathcal{X}_{Gt}$ : grille azimutale de taille $p$ fixée
		\begin{itemize}[label=\textbullet]
			\item $r \in [R_{hub},R]$
			\item $c \in \mathbb{R}$
			\item $\beta \in [0,2\pi]$
			\item $C_L \in \mathcal{C}([0,2\pi],\mathbb{R})$
			\item $C_D \in \mathcal{C}([0,2\pi],\mathbb{R})$
		\end{itemize}
	};
	
	\node[bubble,fill=green!15, right=of XGt] (YGt) { 
		$\mathcal{Y}_{Gt} = \mathbb{R}^{2p}$
		\begin{itemize}[label=\textbullet]
			\item $F_n \in \mathbb{R}^p$
			\item $F_t \in \mathbb{R}^p$
		\end{itemize}
	};
	
	\draw[arrow] (XGt) -- (YGt);
	
\end{tikzpicture}
\caption{Schéma des variantes du modèle global selon la définition de $\mathcal{X}_G$ et $\mathcal{Y}_G$.}
\label{fig model_global_variantes}
\end{figure}
Dans ce formalisme, les caractéristiques locales deviennent des vecteurs de taille $n$. L'espace $\mathcal{X}_{glob}$ contient 1 variable discrète ($B$), 8 scalaires ($\theta, R, R_{hub}, \theta_{pitch}, \rho, \Omega, U_\infty, \gamma$), 2 vecteurs de taille $n$ ($c$ et $\beta$), et les polaires générant deux matrices de taille $n \times m$. \\

La dimension totale de l'espace d'entrée global devient :

$$ d_{glob} = 9 + 2n + 2nm. $$

Le modèle cartographie $\mathcal{X}_{glob} \to \mathbb{R}^{2n}$ pour sortir simultanément les vecteurs de forces sur toute la pale. Cette approche permettrait au réseau de saisir finement comment la géométrie d'une section influence l'aérodynamique des sections adjacentes. \\

Dans la suite, la notation générique $\mathcal{X}$ pourra désigner l'espace d'entrée pour l'approche locale ou l'espace d'entrée pour l'approche globale.



\subsection{Modèle de Correction Résiduelle : La BEM comme compresseur}

Plutôt que d'apprendre la physique de zéro, on propose une architecture qui utilise la BEM classique comme un compresseur d'informations. La BEM calcule une première approximation des efforts $(F_n^{BEM}, F_t^{BEM})$. Le réseau de neurones que nous nommerons \textbf{SkyNet} prend ensuite ces efforts comme base de travail pour corriger les résidus. \\

Pour éviter de surcharger le réseau, nous ne lui redonnons pas l'intégralité de l'espace $\mathcal{X}$, mais un sous-espace réduit $\mathcal{X}'$ (de petite dimension) contenant uniquement les paramètres critiques responsables des écarts entre BEM et Vortex qu'il nous faudra donc déterminer. \\

Ce modèle générique, illustré à la \Cref{fig residual_model}, peut se décliner aussi bien en approche locale qu'en approche globale.

\begin{figure}[H]
    \centering
    \begin{tikzpicture}[
        node distance=1.5cm and 1cm,
        bubble/.style={draw, rectangle, rounded corners=10pt, align=center, fill=blue!10, inner sep=8pt, minimum height=1cm},
        block/.style={draw, rectangle, fill=orange!20, minimum width=2cm, minimum height=1cm, font=\bfseries},
        skynet/.style={draw, rectangle, fill=red!20, minimum width=2.5cm, minimum height=1cm, font=\bfseries\large, rounded corners=3pt},
        arrow/.style={-{Stealth[scale=1.2]}, line width=1pt, draw=gray!80}
    ]
    
    % Noeuds
    \node[bubble] (X) {$x \in \mathcal{X}$};
    \node[block, right=of X] (BEM) {BEM};
    \node[bubble, right=of BEM] (Ybem) {$(F_n^{BEM}, F_t^{BEM}) \in \mathcal{Y}$};
    \node[skynet, right=of Ybem] (SkyNet) {SkyNet};
    \node[bubble, below=of SkyNet] (Yhat) {$(\hat{F}^{Vortex}_n-F^{BEM}_n, \hat{F}^{Vortex}_t-F^{BEM}_t) \in \hat{\mathcal{Y}}_{res}$};
    \node[bubble, fill=green!10, right=of SkyNet] (Xprime) {$x' \in \mathcal{X}'$};
    
    % Flèches
    \draw[arrow] (X) -- (BEM);
    \draw[arrow] (BEM) -- (Ybem);
    \draw[arrow] (Ybem) -- (SkyNet);
    \draw[arrow] (SkyNet) -- (Yhat);
    \draw[arrow] (Xprime) -- (SkyNet);
    
    \end{tikzpicture}
    \caption{Architecture résiduelle où la BEM pré-conditionne les données pour le réseau SkyNet.}
    \label{fig residual_model}
\end{figure}

\subsection{Modèle à Variables Intermédiaires}

Au lieu de prédire directement les forces, SkyNet peut être entraîné pour prédire la cinématique locale du fluide, c'est-à-dire la \textbf{vitesse effective ($V_{eff}$)} et l'\textbf{angle d'attaque ($\alpha$)}. Ces grandeurs constituent des intermédiaires communs aux méthodes BEM et Vortex.\\

Une fois $\alpha$ et $V_{eff}$ prédits, nous pouvons calculer $F_n$ et $F_t$ avec une simple formule aérodynamique. En effet, on peut projeter les coefficients de portance ($C_L$) et de traînée ($C_D$) dans le repère de la section pour obtenir les coefficients normal ($C_n$) et tangentiel ($C_t$) :
\begin{itemize}
    \item $C_n = C_L \cos(\alpha) + C_D \sin(\alpha)$,
    \item $C_t = C_L \sin(\alpha) - C_D \cos(\alpha)$.
\end{itemize}

On transforme ensuite ces coefficients en densités de forces par unité de longueur (N/m) grâce à la pression dynamique :
$$F_n = \frac{1}{2} \rho V_{eff}^2 \cdot c \cdot C_n,$$
$$F_t = \frac{1}{2} \rho V_{eff}^2 \cdot c \cdot C_t.$$

Cette méthode, illustrée à la \Cref{fig kinematic_model}, garantirait que les non-linéarités critiques (comme le décrochage aérodynamique capturé par les tables $C_L/C_D$) sont traitées par la physique exacte et non approximées par le réseau.

\begin{figure}[H]
    \centering
    \begin{tikzpicture}[
        node distance=1.5cm and 1cm,
        bubble/.style={draw, rectangle, rounded corners=10pt, align=center, fill=blue!10, inner sep=8pt, minimum height=1cm},
        block/.style={draw, rectangle, fill=orange!20, minimum width=2.5cm, minimum height=1cm, font=\bfseries, align=center},
        skynet/.style={draw, rectangle, fill=red!20, minimum width=2.5cm, minimum height=1cm, font=\bfseries\large, rounded corners=3pt},
        arrow/.style={-{Stealth[scale=1.2]}, line width=1pt, draw=gray!80}
    ]
    
    % Noeuds
    \node[bubble] (X) {$x \in \mathcal{X}$};
    \node[skynet, right=of X] (SkyNet) {SkyNet};
    \node[bubble, fill=purple!10, right=of SkyNet] (Inter) {$(\hat{\alpha}, \hat{V}_{eff})$};
    \node[block, right=of Inter, text width=3cm] (Eq) {Formules\\Aérodynamiques};
    \node[bubble, right=of Eq] (Y) {$(\hat{F}_n, \hat{F}_t) \in \mathcal{Y}$};
    
    % Flèches
    \draw[arrow] (X) -- (SkyNet);
    \draw[arrow] (SkyNet) -- (Inter);
    \draw[arrow] (Inter) -- (Eq);
    \draw[arrow] (Eq) -- (Y);
    
    \end{tikzpicture}
    \caption{Architecture utilisant des variables intermédiaires avant d'appliquer une formule aérodynamique.}
    \label{fig kinematic_model}
\end{figure}

Cette stratégie générique peut se décliner aussi bien en approche locale qu'en approche globale et elle peut également être couplée à la stratégie de correction résiduelle pour la BEM. \\

En effet, la jonction entre la vision macroscopique de la BEM et la caractérisation locale du profil repose sur la reconstruction du triangle des vitesses au droit de la section. Dans le formalisme BEM, les facteurs d'induction axiale ($a$) et tangentielle ($a'$) modulent respectivement la vitesse incidente $U_\infty$ et la vitesse d'entraînement $\Omega r$ pour définir l'angle de flux $\phi$ via la relation :
\begin{equation}
    \tan \phi = \frac{U_\infty(1-a)}{\Omega r(1+a')}.
\end{equation}

Cette cinématique se traduit directement pour le profil par la vitesse effective $V_{eff}$ et l'angle d'attaque local $\alpha$ :
\begin{equation}
    V_{eff} = \sqrt{\left[ U_\infty(1-a) \right]^2 + \left[ \Omega r (1+a') \right]^2} = \frac{U_\infty(1-a)}{\sin \phi},
\end{equation}
\begin{equation}
    \alpha = \phi - (\beta + \theta_{pitch}).
\end{equation}

Dans une approche Vortex, bien que l'induction soit calculée par l'application de la loi de Biot et Savart à partir de la géométrie du sillage, elle aboutit à la même définition locale du couple $(\alpha, V_{eff})$ pour l'évaluation des forces via les polaires.

\subsection{Recapitualtif}

Il existe donc 8 stratégies possibles :
$$ \text{(local, global)} \times \text{(Avec BEM, Sans BEM)} \times \text{(Variable intermédiaire, Directe)}.$$

\section{Application à l'éolienne New Mexico}

Dans le cadre de l'expérience New Mexico, la plupart des paramètres sont fixés : 
$B=3$, $R_{hub}=\SI{0.210}{\meter}$, $R=\SI{2.25}{\meter}$, $\theta_{pitch}=\SI{-2.30}{\degree}$, $\rho=\SI{1.198}{\kilogram\per\meter\cubed} $, $\Omega=\SI{44.5}{\radian\per\second}$. \\

La grille de discrétisation des rayons est fixée et comporte $n=36$ points. De même, la grille de discrétisation des angles d'attaque est fixée et compte $m=721$ points. Les valeurs des cordes, des twist et des $C_L$ et $C_D$ sont fixées sur ces grilles. \\

En ce qui concerne la discrétisation des azimuts elle est également fixée à $36$ points.

On peut faire varier : la position radiale $r$, la position azimutale $\theta$, le yaw $\gamma$, la vitesse incidente $U_\infty$ (et donc le TSR).

\end{document}